% Contents-only layout, independent of regenerated Markdown content.
\setlength{\cftbeforechapskip}{8pt}
\bookmarksetup{open=true,openlevel=1}
\makeatletter
\newcount\manual@toclevel
\def\manual@splitnumber#1 #2\manual@stop{%
  \def\manual@tocnumber{#1}%
  \def\manual@toctitle{#2}%
}
\newcommand{\manual@numberedcontentsline}[4]{%
  \def\manual@kind{#1}%
  \def\manual@chapter{chapter}%
  \def\manual@section{section}%
  \ifx\manual@kind\manual@chapter
    \manual@toclevel=0
  \else\ifx\manual@kind\manual@section
    \manual@toclevel=1
  \else
    \manual@toclevel=2
  \fi\fi
  \begingroup
    \manual@splitnumber#2\manual@stop
    \ifcase\manual@toclevel
      \settowidth{\cftchapnumwidth}{\cftchapfont\manual@tocnumber\enspace}%
    \or
      \settowidth{\cftsecnumwidth}{\cftsecfont\manual@tocnumber\enspace}%
    \else
      \settowidth{\cftsubsecnumwidth}{\cftsubsecfont\manual@tocnumber\enspace}%
    \fi
    \ifnum\manual@toclevel=2
      \color{tocmuted}\hypersetup{linkcolor=tocmuted}%
    \else
      \color{black}\hypersetup{linkcolor=black}%
    \fi
    \manual@savedcontentsline{#1}{\numberline{\manual@entrystrut\manual@tocnumber}\manual@toctitle}{#3}{#4}%
  \endgroup
}
% Choose column breaks locally; a third-level entry may start a column.
\newcount\manual@total
\newcount\manual@start
\newcount\manual@index
\newcount\manual@natural
\newcount\manual@break
\newcount\manual@candidate
\newcount\manual@prior
\newcount\manual@rank
\newcount\manual@rows
\newcount\manual@shift
\newcount\manual@column
\newdimen\manual@capacity
\newdimen\manual@used
\newdimen\manual@fullheight
\newdimen\manual@remaining
\newdimen\manual@needed
\newdimen\manual@before
\newdimen\manual@strutheight
\newdimen\manual@strutdepth
\newif\ifmanual@legal
\newcommand{\manual@entrystrut}{%
  \vrule width0pt height\manual@strutheight depth\manual@strutdepth}
\def\manual@level#1{\csname manual@level@\number#1\endcsname}
\def\manual@lines#1{\csname manual@lines@\number#1\endcsname}
\def\manual@height#1{\dimexpr
  \ht\csname manual@box@\number#1\endcsname+
  \dp\csname manual@box@\number#1\endcsname\relax}
\def\manual@toc@chapter{0}
\def\manual@toc@section{1}
\def\manual@toc@subsection{2}
\newcommand{\manual@contentsline}[4]{%
  \ifcsname manual@toc@#1\endcsname
    \advance\manual@total\@ne
    \expandafter\newbox\csname manual@box@\the\manual@total\endcsname
    \expandafter\edef\csname manual@level@\the\manual@total\endcsname{%
      \csname manual@toc@#1\endcsname}%
    \setbox\csname manual@box@\the\manual@total\endcsname=\vbox{%
      \hsize=\linewidth \parskip=0pt
      \manual@numberedcontentsline{#1}{#2}{#3}{#4}%
      \expandafter\xdef\csname manual@lines@\the\manual@total\endcsname{%
        \ifnum\prevgraf>0 \the\prevgraf\else 1\fi}%
    }%
    % Give the last line the same minimum depth as an ordinary single-spaced row.
    \ifdim\dp\csname manual@box@\the\manual@total\endcsname<\manual@strutdepth
      \dp\csname manual@box@\the\manual@total\endcsname=\manual@strutdepth
    \fi
  \fi
}
\newcommand{\manual@checkbreak}[1]{%
  \manual@legaltrue
  \manual@prior=#1\relax \advance\manual@prior\m@ne
  \ifnum\manual@prior<\manual@start\else
    % A chapter stays with its next entry; a section with its first subsection.
    \ifnum\manual@level\manual@prior=0
      \manual@legalfalse
    \else\ifnum\manual@level\manual@prior=1
      \ifnum\manual@level#1=2 \manual@legalfalse\fi
    \fi\fi
  \fi
}
\newcommand{\manual@validbreak}{%
  \manual@checkbreak\manual@break
  \ifmanual@legal\else
    \advance\manual@break\m@ne
    \expandafter\manual@validbreak
  \fi
}
\newcommand{\manual@preferheading}{%
  \manual@candidate=\manual@break
  \manual@before=\manual@used
  \manual@rank=3 \manual@rows=0 \manual@shift=0
  \@whilenum\manual@candidate>\manual@start\do{%
    \manual@checkbreak\manual@candidate
    \ifdim\manual@before<\manual@needed \manual@legalfalse\fi
    \ifmanual@legal
      \ifnum\manual@level\manual@candidate=0
        \ifnum\manual@rank>0
          \manual@break=\manual@candidate \manual@rank=0 \manual@shift=\manual@rows
        \fi
      \else\ifnum\manual@level\manual@candidate=1
        \ifnum\manual@rows<3
          \ifnum\manual@rank>1
            \manual@break=\manual@candidate \manual@rank=1 \manual@shift=\manual@rows
          \fi
        \fi
      \fi\fi
    \fi
    \advance\manual@candidate\m@ne
    \advance\manual@rows\manual@lines\manual@candidate\relax
    \advance\manual@before-\manual@height\manual@candidate
    \ifnum\manual@rows>3 \manual@candidate=\manual@start\fi
  }%
}
\newcommand{\manual@fillcolumn}{%
  \manual@used=0pt \manual@index=\manual@start
  \manual@natural=\numexpr\manual@start-1\relax
  \@whilenum\manual@index<\numexpr\manual@total+1\relax\do{%
    \ifdim\dimexpr\manual@used+\manual@height\manual@index\relax>\manual@capacity
      \manual@index=\numexpr\manual@total+1\relax
    \else
      \advance\manual@used\manual@height\manual@index
      \manual@natural=\manual@index \advance\manual@index\@ne
    \fi
  }%
  \manual@break=\numexpr\manual@natural+1\relax
  \ifnum\manual@natural<\manual@total
    \manual@validbreak
    \manual@natural=\numexpr\manual@break-1\relax
    \manual@used=0pt \manual@index=\manual@start
    \@whilenum\manual@index<\manual@break\do{%
      \advance\manual@used\manual@height\manual@index
      \advance\manual@index\@ne
    }%
  \fi
}
\newcommand{\manual@layouttoc}{%
  \manual@start=1 \manual@column=1
  \@whilenum\manual@start<\numexpr\manual@total+1\relax\do{%
    \manual@fullheight=\@colroom
    \ifnum\manual@column>2 \manual@fullheight=\textheight\fi
    \manual@capacity=\manual@fullheight \manual@needed=-\maxdimen
    % Balance the final pair of columns without creating an extra page.
    \ifodd\manual@column
      \manual@remaining=0pt \manual@index=\manual@start
      \@whilenum\manual@index<\numexpr\manual@total+1\relax\do{%
        \advance\manual@remaining\manual@height\manual@index
        \advance\manual@index\@ne
      }%
      \ifdim\manual@remaining>2\manual@fullheight\else
        \manual@needed=\dimexpr\manual@remaining-\manual@fullheight\relax
        \manual@capacity=.5\manual@remaining
        \advance\manual@capacity .5\baselineskip
      \fi
    \fi
    \manual@fillcolumn
    \ifdim\manual@used<\manual@needed
      \manual@capacity=\manual@fullheight \manual@needed=-\maxdimen
      \manual@fillcolumn
    \fi
    \manual@shift=0
    \ifnum\manual@natural<\manual@total
      \manual@preferheading
    \fi
    \typeout{MANUALTOC column=\the\manual@column;
      start=\the\manual@start; natural=\the\manual@natural;
      next=\the\manual@break; shift=\the\manual@shift;
      capacity=\the\manual@capacity}%
    \manual@index=\manual@start
    \@whilenum\manual@index<\manual@break\do{%
      \nointerlineskip
      \box\csname manual@box@\the\manual@index\endcsname
      \penalty\@M
      \advance\manual@index\@ne
    }%
    \manual@start=\manual@break
    \ifnum\manual@start<\numexpr\manual@total+1\relax
      \columnbreak
      \advance\manual@column\@ne
    \fi
  }%
}
\renewcommand{\manualtoc}{%
  \begingroup
    \setstretch{1}\setlength{\parskip}{0pt}%
    \manual@strutheight=\ht\strutbox \manual@strutdepth=\dp\strutbox
    \manual@total=0
    \let\manual@savedcontentsline\contentsline
    \let\contentsline\manual@contentsline
    \raggedcolumns
    \begingroup
      \makeatletter
      \@input{\jobname.toc}%
      \manual@layouttoc
      \if@filesw
        \expandafter\newwrite\csname tf@toc\endcsname
        \immediate\openout\csname tf@toc\endcsname\jobname.toc\relax
      \fi
      \@nobreakfalse
    \endgroup
  \endgroup
}
\makeatother
